\chapter{Podsumowanie i wnioski}

\section{Synteza przeprowadzonych badań}

Praca dotyczyła wykorzystania modelu Qwen3-0.6B do wykrywania niechcianej
poczty na podstawie tematu i treści wiadomości. Publiczne korpusy sprowadzono
do wspólnego schematu, usunięto puste rekordy i duplikaty, a dane treningowe
zbalansowano. Trzy źródła pozostawiono poza treningiem, aby oceniać model także
na wiadomościach o innej charakterystyce.

Model bez dostrajania przypisywał wszystkie oceniane wiadomości do klasy
\texttt{ham}. Na zbalansowanym \texttt{train\_subset} oznaczało to 50,0\%
\emph{accuracy} i F1 równe 0 dla klasy pozytywnej. Po dostrojeniu najlepszy
wynik uzyskała konfiguracja \texttt{K08}: kontekst 512 tokenów, współczynnik
uczenia \(10^{-4}\), LoRA \(r=32\), \(\alpha=64\) oraz dropout 0,05. Siódmy
oceniany punkt kontrolny osiągnął \(S_{\mathrm{ext}}=98{,}0\%\), F1 równe
97,2\% na \texttt{spam\_ham}, \emph{specificity} 96,9\% na \texttt{enron} i
\emph{recall} 99,9\% na \texttt{fraudulent\_email\_corpus}.

Najlepsze warianty generowania etykiety, oceny następnego tokenu i odpowiedzi
ustrukturyzowanej różniły się wynikiem zbiorczym o około 0,1 punktu
procentowego. Najlepszą metodą bazową był DistilBERT z wynikiem
\(S_{\mathrm{ext}}=92{,}3\%\). TF-IDF z Naive Bayes i SVM osiągnęły
odpowiednio 90,7\% i 90,1\%.

\section{Odpowiedzi na pytania badawcze}

\begin{enumerate}
    \item \textbf{Czy dostrojenie małego modelu językowego za pomocą LoRA
    pozwala poprawnie klasyfikować wiadomości ze zbioru treningowego i ze źródeł
    niewykorzystanych do aktualizacji wag?}

    Sam prompt nie wystarczył, aby model bez dostrajania poprawnie wykonywał
    klasyfikację. Po dostrojeniu z wykorzystaniem adaptera LoRA model uzyskał
    wysokie wyniki na podzbiorze danych treningowych oraz na trzech innych
    korpusach.

    \item \textbf{Który z badanych sposobów użycia modelu językowego daje
    najkorzystniejsze połączenie skuteczności, niezależności decyzji od formatu
    generowanej odpowiedzi i prostoty inferencji?}

    Do dalszego porównania wybrano metodę oceny następnego tokenu. Jej przewaga
    liczbowa nad generowaniem etykiety i odpowiedzi ustrukturyzowanej była
    niewielka. Za wyborem przemawiał brak generowania tekstu: decyzja wynikała
    bezpośrednio z porównania prawdopodobieństw etykiet i nie zależała od
    parsera. Klasyfikacja sekwencji również nie wymagała parsera, ale uzyskała
    niższy wynik zewnętrzny, równy 96,5\%.

    \item \textbf{Jak parametry treningu oraz jego czas trwania wpływają
    na wynik modelu na korpusach niewykorzystanych w treningu?}

    Metryki na danych podobnych do treningowych szybko zbliżały się do 100\%,
    przez co słabo rozróżniały konfiguracje. Warianty ze współczynnikiem uczenia
    \(10^{-6}\) uzyskały niższe wyniki na pozostałych korpusach, a kontekst 1024
    tokenów wydłużał trening bez poprawy najlepszego rezultatu. W grupie
    wariantów ze współczynnikiem \(10^{-4}\) najlepiej wypadł adapter z
    \(r=32\) i \(\alpha=64\). Mniejszy adapter uzyskał słabszy wynik, natomiast
    zwiększenie pojemności ponad \(r=32\) i \(\alpha=64\) nie dawało regularnej
    poprawy. Nie udało się określić jednoznacznego wpływu dropout. Dla
    najlepszych konfiguracji najwyższe \(S_{\mathrm{ext}}\) pojawiało się przed
    końcem treningu. Wybór
    ostatniego adaptera lub kierowanie się wyłącznie metrykami treningowymi
    prowadziłoby do wyboru słabszego wariantu.

    \item \textbf{Jaki kompromis między skutecznością a kosztem inferencji
    wynika z porównania dostrojonego modelu językowego z prostszymi
    klasyfikatorami?}

    Qwen3 uzyskał wynik \(S_{\mathrm{ext}}\) wyższy o 5,7 punktu procentowego
    od DistilBERT, ale był od niego około ośmiokrotnie wolniejszy. Metody TF-IDF
    przetwarzały od 2,2 do 10,6 tys. wiadomości na sekundę, przy wyniku
    zewnętrznym niższym o 7,3--18,4 punktu procentowego. Przepustowość Qwen3,
    równa 7,3 wiadomości na sekundę, może ograniczać zastosowanie tego modelu
    jako jedynego filtra całego ruchu. Pomiary wykonane na jednym komputerze
    nie wystarczają jednak do wyceny konkretnego wdrożenia.
\end{enumerate}

\section{Ograniczenia pracy}

Eksperymenty przeprowadzono na publicznych korpusach, z których część powstała
wiele lat temu. Wiadomości z poszczególnych źródeł różnią się formatowaniem,
słownictwem i sposobem nadawania etykiet. Deduplikacja ograniczyła liczbę
powtórzeń, ale nie usunęła wszystkich cech pozwalających rozpoznać źródło
wiadomości. Wyników nie potwierdzono na bieżącym strumieniu poczty.

Trzy zbiory zewnętrzne nie były używane do treningu Qwen3, ale ich wyniki
służyły do wyboru konfiguracji i punktu kontrolnego. W metodach bazowych
parametry dobierano wyłącznie na części walidacyjnej, a wymienione zbiory
wykorzystano dopiero do oceny wybranego wariantu.

Każdy trening wykonano z jednym ziarnem losowym. Nie obliczano przedziałów
ufności i nie powtarzano tych samych konfiguracji, dlatego różnice rzędu
0,1--0,3 punktu procentowego mogą wynikać częściowo z losowości treningu lub
próbkowania. Dotyczy to zwłaszcza kolejności najlepszych sposobów dostrajania.

Zakres strojenia nie był jednakowy. Szesnaście konfiguracji i pośrednie punkty
kontrolne oceniono dla metody następnego tokenu. Pozostałe metody korzystały z
wybranej części tej przestrzeni, a dla odpowiedzi ustrukturyzowanej ukończono
dwie konfiguracje. Dla metod bazowych sprawdzono ograniczone zbiory parametrów
i jeden podział walidacyjny. Naive Bayes nie podlegał strojeniu, a DistilBERT
wytrenowano tylko z jednym zestawem parametrów.
Przedstawione wyniki dotyczą więc wyłącznie sprawdzonych konfiguracji.

Badanie objęło jeden model bazowy i jeden jego rozmiar. Pomiar wydajności
wykonano na jednym komputerze i dla jednej próbki. Nie sprawdzono zachowania
przy wielu równoczesnych żądaniach, opóźnień krańcowych, zużycia energii ani
kosztu pracy serwera. Zmierzone wartości pokazują względną różnicę między
metodami w badanym środowisku, ale nie opisują infrastruktury produkcyjnej.

Klasyfikacja binarna łączy spam reklamowy, phishing i wiadomości oszukańcze w
jedną klasę. Model nie określa rodzaju zagrożenia ani nie wyjaśnia, który
fragment wiadomości wpłynął na decyzję. Do modelu przekazywano tylko temat i
treść, bez reputacji nadawcy, pełnych nagłówków i mechanizmów uwierzytelniania
domeny.

\section{Możliwość wdrożenia}

Koszt inferencji porównano na wspólnej próbce 1000 wiadomości. Qwen3 z
adapterem LoRA przetwarzał 7,3 wiadomości na sekundę, a przetworzenie całej
próbki trwało 136,7 sekundy. Proces wykorzystał szczytowo 3466~MiB pamięci
operacyjnej, natomiast model bazowy i adapter zajmowały razem około 1526~MiB na
dysku. DistilBERT przetwarzał 57,9 wiadomości na sekundę, wykorzystywał
szczytowo 563~MiB pamięci i zajmował 256~MiB na dysku. Dla TF-IDF z Naive
Bayes przepustowość wyniosła około 10,6 tys. wiadomości na sekundę, szczytowe
użycie pamięci 337~MiB, a rozmiar modelu z wektoryzatorem 13~MiB.

Zmierzona przepustowość Qwen3 ogranicza możliwość wykorzystania go jako
jedynego filtra pierwszej linii bez skalowania infrastruktury. Szybszy sprzęt,
kwantyzacja i przetwarzanie wiadomości w partiach mogłyby poprawić
przepustowość. Na podstawie wykonanych
testów nie można oszacować kosztu konkretnego wdrożenia serwerowego.

DistilBERT osiągnął najlepszy wynik wśród metod bazowych i działał szybciej od
Qwen3, ale nadal był wielokrotnie wolniejszy od TF-IDF.

Koszt inferencji można byłoby ograniczyć przez zastosowanie układu kaskadowego.
Uwierzytelnianie nadawcy, reguły, reputacja i tani klasyfikator tekstu
obsługiwałyby przypadki jednoznaczne, a model językowy analizowałby tylko
pozostałe wiadomości. W pracy nie sprawdzono, czy Qwen3 poprawia wyniki właśnie
dla przykładów, co do których prostszy klasyfikator jest niepewny.

Model Qwen3-0.6B z adapterem LoRA jest prototypem klasyfikatora treści, a nie
kompletnym filtrem pocztowym. W systemie produkcyjnym jego wynik mógłby być
jednym z sygnałów używanych do podjęcia decyzji. Samodzielne odrzucanie
wiadomości wymagałoby dokładniejszej kontroli fałszywych alarmów i testów na
danych produkcyjnych.

\section{Kierunki dalszych badań}

Wyniki odpowiedzi ustrukturyzowanej trzeba uzupełnić o dwie brakujące
konfiguracje. Szersze przeszukanie parametrów Naive Bayes i DistilBERT
pozwoliłoby sprawdzić, czy ich obecne wyniki są bliskie najlepszym dla przyjętych
danych. Osobno trzeba dobrać progi klasyfikacji do założonego kosztu
\emph{false positive} i \emph{false negative}.

Kolejny eksperyment powinien rozdzielić wybór konfiguracji Qwen3 od testu na
zbiorach zewnętrznych. Można również przeprowadzić walidację
\emph{leave-one-source-out}, kolejno wyłączając z treningu każde źródło. Innym
wariantem jest podział czasowy, w którym do testu trafiają nowsze wiadomości.
Pomiar na niezbalansowanym strumieniu pozwoliłby
ocenić liczbę fałszywych alarmów przy proporcjach klas zbliżonych do ruchu
produkcyjnego. Osobne etykiety
spamu reklamowego, phishingu i oszustw pozwoliłyby sprawdzić, czy model
rozpoznaje także rodzaj zagrożenia.

Ocena wdrożeniowa wymaga porównania wariantów kwantyzacji, silników inferencji
i ustawień przetwarzania partiami na sprzęcie serwerowym. Oprócz średniej
przepustowości należy zmierzyć opóźnienie, użycie pamięci przy równoczesnych
żądaniach oraz koszt przetworzenia określonej liczby wiadomości. Osobnego testu
wymaga też układ kaskadowy z tanim klasyfikatorem pierwszego etapu.

W badanych warunkach dostrojony Qwen3-0.6B uzyskał wyższy wynik zbiorczy niż
metody bazowe, lecz wymagał znacznie więcej czasu i pamięci. Qwen3 mógłby
pełnić rolę klasyfikatora drugiego etapu, ale takiego układu nie oceniono w
eksperymentach.
